\chapter{RECOMENDACIONES}

Con base en las experiencias metodológicas, la evaluación experimental y los hallazgos alcanzados en la presente investigación, se formulan las siguientes recomendaciones:

\noindent \textbf{PRIMERA:} Se recomienda trasladar el modelo de clasificación híbrido evaluado en entorno de laboratorio hacia infraestructuras de prueba en campo en microempresas (MYPEs/PYMEs) del departamento de Puno, desplegando los artefactos optimizados en formato de intercambio abierto ONNX (\textit{Open Neural Network Exchange}) sobre pasarelas IoT físicas de bajo costo (tales como Raspberry Pi 4/5, NVIDIA Jetson Orin Nano o pasarelas industriales x86/ARM64). Esta proyección técnica permitirá validar el rendimiento del sistema ante variaciones dinámicas del tráfico en tiempo real y condiciones heterogéneas de conectividad inalámbrica.

\vspace{0.8em}
\noindent \textbf{SEGUNDA:} Se recomienda ampliar y precisar los criterios de recolección de datos incorporando mecanismos de captura continua de flujos de red en tiempo real mediante zócalos (\textit{sockets}) y bibliotecas de inspección profunda de paquetes (tales como \texttt{pcap} o \texttt{DPDK}). Asimismo, se sugiere explorar esquemas adaptativos de sobremuestreo sintético con sintonización dinámica del umbral de frontera en Borderline-SMOTE1, combinados con técnicas de selección de características de mayor orden espacial para fortalecer la calidad de los datos de entrada ante variaciones del tráfico.

\vspace{0.8em}
\noindent \textbf{TERCERA:} Se recomienda realizar comparaciones entre diferentes meta-modelos de Nivel 1 y variaciones en la cuantificación de incertidumbre, explorando clasificadores basados en CatBoost GPU, TabNet o redes profundas tabulares, así como formulaciones entrópicas alternativas (tales como la Entropía de Tsallis o Rényi). Asimismo, se sugiere incorporar módulos de explicabilidad algorítmica ágil mediante TreeSHAP, permitiendo a los analistas de Centros de Operaciones de Seguridad (SOC) auditar visualmente las razones por las cuales una conexión de red fue clasificada como ataque minoritario.

\vspace{0.8em}
\noindent \textbf{CUARTA:} Se recomienda establecer alianzas y protocolos de validación en colaboración con especialistas en ciberseguridad industrial y administradores de redes de pequeñas y medianas empresas. Esto permitirá contrastar las alertas automáticas generadas por el Stacking Híbrido con logs de auditoría física, asegurando que la solución no solo alcance métricas cuantitativas sobresalientes ($F1\text{-}Macro = 70.18\%$, $\mathrm{Latencia} = 8.65\,\mu s$), sino que mantenga congruencia operativa con los requerimientos de respuesta ante incidentes cibernéticos en entornos reales.